\documentclass[10pt,a4paper]{amsart}
\usepackage[margin=19mm]{geometry}
\usepackage{amsmath,amssymb,mathtools}
\usepackage[scr=rsfs,cal=euler]{mathalfa}
\usepackage{xfrac}
\usepackage[unicode,hidelinks,bookmarks=false]{hyperref}
\newcommand{\ie}{i.e.\ }
\newcommand{\cf}{cf.\ }
\newcommand{\resp}{respectively\ }
\newcommand{\Subsection}{\subsection}
\providecommand{\nobreakdash}{\nobreak\hbox{-}}
\setlength{\emergencystretch}{2em}
\multlinegap0pt
\usepackage{bm}
\usepackage{bbm}
\usepackage{dsfont}
\usepackage{xspace}
\usepackage{cancel}
\usepackage{mathtools}
\usepackage{amsthm}
\makeatletter
\@ifundefined{theorem}{\newtheorem{theorem}{Theorem}}{}
\@ifundefined{prop}{\newtheorem{prop}[theorem]{Proposition}}{}
\makeatother
\makeatletter
\newcommand{\Z}{\mathbb{Z}}
\expandafter\ifx\csname abstract\endcsname\relax
\renewenvironment{abstract}{
  \vspace{1em}
  \begin{list}{}{
    \setlength{\leftmargin}{1cm} 
    \setlength{\rightmargin}{1cm}
  }
  \item[]
  \noindent{ \scshape Abstract.}\hspace{0.3em}\footnotesize
}{
  \end{list}
  \par\vspace{1em}
}
\fi
\makeatother
\title[Companion matrices, permutations and lattice ideals]{Companion matrices, permutations and lattice ideals}
\author{Nsibiet E. Udo, Praise Adeyemo}
\date{Source: arXiv:2511.06555v1 (CC BY 4.0)}
\begin{document}
\maketitle

\noindent\emph{Source and licence.} Companion matrices, permutations and lattice ideals. Version arXiv:2511.06555v1 (CC BY 4.0), distributed under the Creative Commons Attribution 4.0 International licence, as recorded in the arXiv licence metadata for this exact version and shown on the abstract page https://arxiv.org/abs/2511.06555v1. Full rights notice: licenses/arxiv-2511.06555-CC-BY-4.0.txt.\par\smallskip

\noindent\textit{Editorial scope.} Counted unit: the Prop whose source label is \texttt{prop:isomorphism-btw-Gn-quotient-Ln} in the article (source0.tex lines 627--630, source label \texttt{prop:isomorphism-btw-Gn-quotient-Ln}), delivered together with the article's own complete original proof of it (source0.tex lines 632--712, one \texttt{proof} environment of 81 lines). No mathematics is written, repaired or re-derived: statement and proof are the article's own text line for line. Required context retained uncounted: none. Every definition, display and label of those lines is retained unchanged. Omitted from this package and not redistributed: 1--626, 631--631, 713--1021. Nothing inside a retained excerpt is omitted or altered: every retained body is delivered exactly as the article's lines are written, so no omission receipt of any kind accompanies this package. Citations: the retained excerpts carry no citation command at all, so no bibliography entry of the article is transcribed and no reference is redistributed. Reference adaptation: none was needed and none was made; every reference inside the delivered unit resolves inside it, so no unresolved reference or citation is produced. Printed slips kept literal: no printed word, symbol, hypothesis or number of the counted statement or of its proof was altered. Layout and class: the article's own class is \texttt{amsart} with packages including fontenc, lmodern, babel, geometry, amssymb, amsthm, verbatim, latexsym, amsfonts, lscape, amsmath, graphicx, youngtab, rotating; the wrapper is the lane's pinned \texttt{amsart} editorial wrapper at 10pt on A4, which restates the theorem-like environments the retained text needs and adds the article's own macro definitions that the retained text needs (\texttt{\textbackslash Z}), with extra wrapper packages (bm, bbm, dsfont, xspace, cancel, mathtools). The delivered numbering is the unit's own. Assets: no retained span contains a figure environment, a graphics-inclusion command, a PDF-page inclusion, a table or a TikZ picture; no image file is copied into this package, the package declares no asset dependency and no image file is redistributed with it. Licence: version arXiv:2511.06555v1 of this article is distributed under the Creative Commons Attribution 4.0 International licence, as recorded in the arXiv licence metadata for this exact version and shown on the abstract page https://arxiv.org/abs/2511.06555v1. A full rights notice is delivered at licenses/arxiv-2511.06555-CC-BY-4.0.txt. Characters: every retained excerpt is pure ASCII, so the delivered engine, XeLaTeX with its default Latin Modern text font, reports no missing character.
\begin{prop}
	\label{prop:isomorphism-btw-Gn-quotient-Ln}
	The group \( G_n \) is isomorphic to \( \mathbb{Z}^n / L_n \).
\end{prop}

\begin{proof}
	The lattice \( L_n \subset \mathbb{Z}^n \) is generated by the vectors
	\begin{equation}
	    \begin{cases}
	v_i = 2e_i - \sum_{j=1}^n e_j & (1 \leq i \leq n), \\
	v_{n+1} = (2n - 4)e_n, \\
	v_{i+n+1} = 2e_i - 2H(n-3)e_n & (1 \leq i \leq n-1),
	\end{cases} 
    \label{eqn:s5n1}
	\end{equation}
	where \( H(n-3) \) is the Heaviside step function, \( e_1, \dots, e_n\) are canonical basis for \( \Z^n \). Since \( L_n \) has rank \( n \), we may select \( n \) linearly independent generators, such as \( \{v_1, \dots, v_n\} \). These form the columns of the relation matrix
	\[
	R := 2\boldsymbol{I}_n - Q_n \in M_n(\Z),
	\]
	where \( Q_n \) is the all-ones matrix, and \( \boldsymbol{I}_n\) the \( n \times n \) identity. Explicitly,
	\[
	R = 
	\begin{pmatrix}
	1 & -1 & -1 & \cdots & -1 \\
	-1 & 1 & -1 & \cdots & -1 \\
	-1 & -1 & 1 & \cdots & -1 \\
	\vdots & \vdots & \vdots & \ddots & \vdots \\
	-1 & -1 & -1 & \cdots & 1
	\end{pmatrix}.
	\]
	
	The group is isomorphic to the cokernel of the linear map \( R: \mathbb{Z}^n \to \mathbb{Z}^n \), i.e., \( G_n \cong \mathbb{Z}^n / \operatorname{im}(R) \). To determine the structure of \( \mathbb{Z}^n / L_n \), we compute the Smith normal form (SNF) of \( R \). The invariant factors of \( R \) will reveal the group structure and the index \( [\mathbb{Z}^n : L_n] = |\det(R)| \).
	
	Define the all-ones vector \( \mathbf{1} \) and the vector \( f \) as
	\[
	\mathbf{1} = \sum_{i=1}^n e_i = (1,1,\dots,1)^T, \quad f = \sum_{i=2}^n e_i = (0,1,1,\dots,1)^T \in \mathbb{Z}^n,
	\]
	and consider the unimodular matrices
	\[
	U = \boldsymbol{I}_n + f e_1^T, \quad V = \boldsymbol{I}_n + e_1 f^T.
	\]
	These matrices correspond to sequences of elementary row and column operations: subtracting the first row/column from the others, then adding twice the first row/column to the others. One may verify that \( \det(U) = \det(V) = 1 \), since \( e_1^T f = 0 \).
	
	Now compute the product
	\[
	\begin{aligned}
	URV &= (I + f e_1^T)(2I - \mathbf{1}\mathbf{1}^T)(I + e_1 f^T) \\
	&= (I + f e_1^T)\left[2(I + e_1 f^T) - \mathbf{1}\mathbf{1}^T(I + e_1 f^T)\right] \\
	&= (I + f e_1^T)\left[2I + 2e_1 f^T - \mathbf{1}\mathbf{1}^T - \mathbf{1}(e_1^T \mathbf{1}) f^T\right] \\
	&= (I + f e_1^T)\left[2I + (2e_1 - \mathbf{1})f^T - \mathbf{1}\mathbf{1}^T\right].
	\end{aligned}
	\]
	Multiplying out and simplifying using the identities \( e_1^T f = 0 \), \( e_1^T \mathbf{1} = 1 \), and \( f^T \mathbf{1} = n-1 \), we obtain
	\[
	URV = 2I - e_1 e_1^T - 2f f^T = \begin{pmatrix} 1 & \mathbf{0}^T \\ \mathbf{0} & M' \end{pmatrix},
	\]
	where \( M' = -2(Q_{n-1} - \boldsymbol{I}_{n-1}) \).
	
	Now, the matrix \( K = Q_{n-1} - \boldsymbol{I}_{n-1} \) has eigenvalues \( n-2 \) (with eigenvector \( \mathbf{1} \)) and \( -1 \) (with multiplicity \( n-2 \)). Hence,
	\[
	\det(K) = (n-2)(-1)^{n-2}.
	\]
	Since all entries of \( K \) are 0 or 1, the greatest common divisor of its entries is 1, so the first invariant factor of \( K \) is 1. The invariant factors \( d_1 \mid d_2 \mid \cdots \mid d_{n-1} \) must satisfy \( d_1 \cdots d_{n-1} = |\det(K)| = n-2 \), and since \( d_1 = 1 \) and $K$ is an integer matrix of full rank $n-1$, it follows that
	\[
	\operatorname{SNF}(K) = \operatorname{diag}(1, 1, \dots, 1, n-2),
	\]
	with the entry \(1\) occurring \( n-2 \) times. Therefore,
	\[
	\operatorname{SNF}(M') = \operatorname{diag}(2, 2, \dots, 2, 2(n-2)),
	\]
	with \( 2 \) repeated \(n - 2\) times, and combining with the leading 1 gives
	\[
	\operatorname{SNF}(R) = \operatorname{diag}(1, 2, 2, \dots, 2, 2(n-2)).
	\]
	
	Thus,
	\[
	\mathbb{Z}^n / L_n \cong \mathbb Z^n /  \operatorname{im}(R) \cong (\mathbb Z/2 \mathbb Z)^{\,n-2}  \oplus \mathbb{Z}/(2(n-2))\mathbb Z,
	\]
	with \( n-2 \) copies of \( \mathbb{Z}/2 \), and
	\[
	|\mathbb{Z}^n / L_n| = |\det(R)| = 1 \cdot 2^{n-2} \cdot 2(n-2) = (n-2)2^{n-1}.
	\]
	This completes the proof. 
	
\end{proof}

\end{document}
